\documentclass[10pt,a4paper]{amsart}
\usepackage[margin=19mm]{geometry}
\usepackage{amsmath,amssymb,mathtools}
\usepackage[scr=rsfs,cal=euler]{mathalfa}
\usepackage{xfrac}
\usepackage[unicode,hidelinks,bookmarks=false]{hyperref}
\newcommand{\ie}{i.e.\ }
\newcommand{\cf}{cf.\ }
\newcommand{\resp}{respectively\ }
\newcommand{\Subsection}{\subsection}
\providecommand{\nobreakdash}{\nobreak\hbox{-}}
\setlength{\emergencystretch}{2em}
\multlinegap0pt
\usepackage{amssymb}
\newtheorem{theorem}{Theorem}
\title{Petrov type extension for multivalued contraction mappings}
\author{Hakan Sahin, Mustafa Aslantas, Ishak Altun}
\date{Source: arXiv:2605.13358v1 (CC BY 4.0)}
\begin{document}
\maketitle

\noindent\emph{Source and licence.} Petrov type extension for multivalued contraction mappings. Version arXiv:2605.13358v1 (CC BY 4.0), distributed by arXiv under the Creative Commons Attribution 4.0 International licence, http://creativecommons.org/licenses/by/4.0/, as recorded in the arXiv licence metadata for this exact version and shown on the abstract page https://arxiv.org/abs/2605.13358v1. Full licence notice: \texttt{licenses/arxiv-2605.13358-CC-BY-4.0.txt}.\par\smallskip

\noindent\textit{Editorial scope.} Counted unit: the article's theorem Theo1 (source0.tex lines 459--465). The article's complete original proof of it is delivered with it (source0.tex lines 467--608, one \texttt{proof} environment of 142 lines). No mathematics is written, repaired or re-derived: the statement and the proof are the article's own lines, in the article's own notation, and no retained line is substituted or re-flowed.  No further source-local context is needed: every cross-reference of every retained line resolves inside the two delivered spans.  Nothing was omitted from any retained span, so the package carries no omission record.  No retained line contains a citation, so the wrapper carries no bibliography and no bibliography entry.  No retained line contains a figure environment, an image-inclusion command or drawing code, so no image is delivered and the package declares no asset dependency.  Editorial formatting only, in addition: an AMS \texttt{amsart} preamble that restates the article's own declarations and macros used by the retained lines, namely the environments \texttt{theorem} on separate counters, together with the standard packages those lines need.  The retained environments print in this document as Theorem 1 in order of appearance, whereas the article prints the same items with the numbering of its own full document.  The complete native source of the article as distributed in the arXiv e-print for this exact version is bundled unchanged as \texttt{sources/200496/source0.tex}.
\begin{theorem}
\label{Theo1}Let $\left( X,\rho \right) $ be a complete metric space with $%
\left \vert X\right \vert \geq 3$ and $T:X\rightarrow CB(X)$ be a
multivalued $\lambda $-contracting perimeters of triangles on $X$ with $%
\lambda \in \left[ 0,\frac{1}{2}\right) $. If the mapping $T~$has the
property of forming a triangle, then $T$ has a fixed point.
\end{theorem}

\begin{proof}
Except for the obvious situation, we can assume that $\lambda >0$ in the
proof. Let $x_{0}\in X$ and $x_{1}\in Tx_{0}$. If $x_{0}=x_{1}$, then the
proof is finished. Hence, we assume that $x_{0}\neq x_{1}$. Then, since $T$
has the property of forming a triangle, there exists $x_{2}\in Tx_{1}$ with $%
H(Tx_{0},Tx_{2})>0$, and so $x_{0}\neq x_{2}$ such that 
\begin{equation}
\rho (x_{1},x_{2})\leq H(Tx_{0},Tx_{1})+\lambda .  \label{E1}
\end{equation}%
Now, we can consider the condition $x_{1}\neq x_{2}$ (otherwise the proof is
finished). Using the property of forming a triangle of $T$ there exists $%
x_{3}\in Tx_{2}~$with $H(Tx_{1},Tx_{3})>0$, and so $x_{1}\neq x_{3}$ such
that 
\begin{equation}
\rho (x_{2},x_{3})\leq H(Tx_{1},Tx_{2})+\min \left \{ \lambda
^{2},H(Tx_{0},Tx_{2})\right \} .  \label{E2}
\end{equation}

Continuing this process, we can construct a sequence $\left \{
x_{n}\right
\} $ whose three consecutive terms of the sequence are
different from each other with $H\left( Tx_{n-2},Tx_{n}\right) >0$ such that%
\begin{equation*}
\rho (x_{1},x_{2})\leq H(Tx_{0},Tx_{1})+\lambda
\end{equation*}%
and 
\begin{equation}
\rho (x_{n},x_{n+1})\leq H(Tx_{n-1},Tx_{n})+\min \left \{ \lambda
^{n},H(Tx_{n-2},Tx_{n})\right \}  \label{E4}
\end{equation}%
for all $n\geq 2$. Using inequalities $\left( \ref{E1}\right) $ and $\left( %
\ref{E2}\right) $, we have 
\begin{eqnarray}
\rho (x_{1},x_{2})+\rho (x_{2},x_{3}) &\leq &H(Tx_{0},Tx_{1})+\lambda
+H(Tx_{1},Tx_{2})  \notag \\
&&+\min \left \{ \lambda ^{2},H(Tx_{0},Tx_{2})\right \}  \notag \\
&\leq &H(Tx_{0},Tx_{1})+H(Tx_{1},Tx_{2})+H(Tx_{0},Tx_{2})+\lambda  \notag \\
&\leq &\lambda \left( \rho (x_{0},x_{1})+\rho (x_{1},x_{2})+\rho
(x_{0},x_{2})\right) +\lambda  \notag \\
&\leq &2\lambda \left( \rho (x_{0},x_{1})+\rho (x_{1},x_{2})\right) +\lambda
.  \label{E6}
\end{eqnarray}%
In a similar way, from the inequality $\left( \ref{E4}\right) $, we get%
\begin{eqnarray}
\rho (x_{2},x_{3})+\rho (x_{3},x_{4}) &\leq &H(Tx_{1},Tx_{2})+\min \left \{
\lambda ^{2},H(Tx_{0},Tx_{2})\right \}  \notag \\
&&+H(Tx_{2},Tx_{3})+\min \left \{ \lambda ^{3},H(Tx_{1},Tx_{3})\right \} 
\notag \\
&\leq &H(Tx_{1},Tx_{2})+H(Tx_{2},Tx_{3})+H(Tx_{1},Tx_{3})+\lambda ^{2} 
\notag \\
&\leq &\lambda \left( \rho (x_{1},x_{2})+\rho (x_{2},x_{3})+\rho
(x_{1},x_{3})\right) +\lambda ^{2}  \notag \\
&\leq &2\lambda \left( \rho (x_{1},x_{2})+\rho (x_{2},x_{3})\right) +\lambda
^{2}.  \label{E7}
\end{eqnarray}%
Hence, we get 
\begin{equation*}
\rho (x_{n},x_{n+1})+\rho (x_{n+1},x_{n+2})\leq 2\lambda \left( \rho
(x_{n-1},x_{n})+\rho (x_{n},x_{n+1})\right) +\lambda ^{n}
\end{equation*}%
for all $n\geq 1$. Hence, we conclude that 
\begin{eqnarray*}
\rho (x_{n},x_{n+1})+\rho (x_{n+1},x_{n+2}) &\leq &2\lambda \left( \rho
(x_{n-1},x_{n})+\rho (x_{n},x_{n+1})\right) +\lambda ^{n} \\
&\leq &2\lambda \left \{ 2\lambda \left( \rho (x_{n-2},x_{n-1})+\rho
(x_{n-1},x_{n})\right) +\lambda ^{n-1}\right \} \\
&&+2\lambda ^{n} \\
&=&\left( 2\lambda \right) ^{2}\left( \rho (x_{n-2},x_{n-1})+\rho
(x_{n-1},x_{n})\right) +4\lambda ^{n} \\
&\leq &\left( 2\lambda \right) ^{2}\left \{ 2\lambda \left( \rho
(x_{n-3},x_{n-2})+\rho (x_{n-2},x_{n-1})\right) +\lambda ^{n-2}\right \} \\
&&+8\lambda ^{n} \\
&=&\left( 2\lambda \right) ^{3}\left( \rho (x_{n-3},x_{n-2})+\rho
(x_{n-2},x_{n-1})\right) +8\lambda ^{n} \\
&\leq &\left( 2\lambda \right) ^{3}\left \{ 2\lambda \left( \rho
(x_{n-4},x_{n-3})+\rho (x_{n-3},x_{n-2})\right) +\lambda ^{n-3}\right \} \\
&&+8\lambda ^{n} \\
&=&\left( 2\lambda \right) ^{4}\left( \rho (x_{n-4},x_{n-3})+\rho
(x_{n-3},x_{n-2})\right) +16\lambda ^{n} \\
&&\vdots \\
&\leq &\left( 2\lambda \right) ^{n}\left( \rho (x_{0},x_{1})+\rho
(x_{1},x_{2})\right) +2^{n}\lambda ^{n}
\end{eqnarray*}%
for all $n\geq 1$. Hence, we have 
\begin{equation*}
\sum \limits_{n=1}^{\infty }\rho (x_{n},x_{n+1})\leq \sum
\limits_{n=1}^{\infty }\left( 2\lambda \right) ^{n}\left( \rho
(x_{0},x_{1})+\rho (x_{1},x_{2})+1\right) <\infty .
\end{equation*}

So that, $\left \{ x_{n}\right \} $ is a Cauchy sequence. Because of the
completeness of the space $(X,\rho )$, there exist a point $x^{\ast }$ in $X$
such that $x_{n}\rightarrow x^{\ast }$. Let define a set 
\begin{equation*}
A=\{n_{i}:x_{n_{i}}=x^{\ast }\text{ for }i\in 
%TCIMACRO{\U{2115} }%
%BeginExpansion
\mathbb{N}
%EndExpansion
\}.
\end{equation*}%
First, we claim that $A$ is finite set. Assume the contrary. Then, since $%
x_{n}$ is not a fixed point of the mapping $T$ for all $n\in 
%TCIMACRO{\U{2115} }%
%BeginExpansion
\mathbb{N}
%EndExpansion
$, we have 
\begin{equation*}
0<\rho \left( x_{n_{i}},Tx_{n_{i}}\right) =\rho (x^{\ast },Tx^{\ast })
\end{equation*}%
for all $n_{i}\in A$. Also, for all $n_{i}\in A$ we have 
\begin{equation*}
\rho (x^{\ast },Tx^{\ast })=\rho \left( x_{n_{i}},Tx_{n_{i}}\right) \leq
\rho \left( x_{n_{i}},x_{n_{i}+1}\right) \rightarrow 0
\end{equation*}%
as $n_{i}\rightarrow \infty $ which is a contradiction. Hence, $A$ is finite
set. Now, let 
\begin{equation*}
n_{0}=\left \{ 
\begin{array}{ccc}
0 & , & A=\emptyset \\ 
&  &  \\ 
\max A & , & A\neq \emptyset%
\end{array}%
\right. .
\end{equation*}%
Then, we have $x_{n}\neq x^{\ast }$ for all $n>n_{0}$. Hence, we get 
\begin{eqnarray*}
\rho \left( x_{n+1},Tx^{\ast }\right) &\leq &H(Tx_{n},Tx^{\ast
})+H(Tx_{n},Tx_{n+1})+H(Tx_{n+1},Tx^{\ast }) \\
&\leq &\lambda (\rho (x_{n},x^{\ast })+\rho (x_{n},x_{n+1})+\rho
(x_{n+1},x^{\ast }))
\end{eqnarray*}%
for all $n>n_{0}$. Taking limit $n\rightarrow \infty $ in the last
inequality we have 
\begin{equation*}
\rho \left( x^{\ast },Tx^{\ast }\right) =0.
\end{equation*}%
So that, we get $x^{\ast }\in \overline{Tx^{\ast }}=Tx^{\ast }$. That is, $%
x^{\ast }$ is a fixed point of $T$.
\end{proof}

\end{document}
